\documentclass[12pt,a4paper]{article}
% \documentcla

\usepackage{fontspec}

\usepackage[spanish,mexico]{babel}
\usepackage[a4paper,margin=2cm]{geometry}
\usepackage{microtype}
\usepackage{xcolor}

\usepackage{tikz}
\usepackage{array,tabularx,booktabs}
\usepackage{enumitem}
\usetikzlibrary{arrows.meta,positioning,shapes.geometric,fit}


\usepackage{graphicx}

\usepackage{listings}
\usepackage[hidelinks]{hyperref}

\setlength{\parindent}{0pt}

% %
\setlength{\parskip}{5pt}
\setlist{nosep,leftmargin=*}

\renewcommand{\arraystretch}{1.15}

\definecolor{codebg}{RGB}{246,246,246}
\definecolor{codeframe}{RGB}{210,210,210}
\lstdefinestyle{gocode}{
  language=Go,
  basicstyle=\ttfamily\small,
  backgroundcolor=\color{codebg},
  frame=single,
  rulecolor=\color{codeframe},
  breaklines=true,
  columns=fullflexible,
  showstringspaces=false,
  tabsize=4
}
\newcommand{\pendiente}[1]{\textbf{[Completar: #1]}}

\tikzset{
  bloque/.style={draw,rounded corners,align=center,minimum height=9mm,
    text width=3.0cm,inner sep=5pt},
  flecha/.style={-{Latex[length=2mm]},thick},
  estado/.style={draw,rounded corners,align=center,minimum width=2.5cm,
    minimum height=9mm,inner sep=4pt}
}

%%% 

\setlength{\parindent}{20pt}

\begin{document}

\title{\LARGE\textbf{Reporte}}

\author{Universidad Nacional Autónoma de México\\[0.6em]
  Modelado y Programación\\[0.6em]
  Facultad de Ciencias}
\date{}
\maketitle
\begin{center}
\rule{\linewidth}{0.4pt}\\[1.2em]
\textbf{Jafet Felipe Flores}\\[1em]
\rule{\linewidth}{0.4pt}\\[1.2em]
\graphicspath{ {./Images/} }
\includegraphics[scale=.8]{Portada.png}
\end{center}


\newpage

\section*{Resumen}
Este proyecto trata de implementar y diseñar un chat a través de una comunicación
cliente-servidor, el cual será desarrollado en el lenguaje de programación Go,
tanto el cliente como el servidor. La comunicación se basa en un intercambio de
mensajes JSON. El servidor se encargará de procesar las conexiones, que cada
cliente se identifique antes de poder hacer cualquier otra cosa, así también como
responder a los JSON recibidos según sea cada caso. Por la parte del cliente,
es quien manda los mensajes al servidor para ser debidamente procesados. Se hace
uso de la concurrencia para poder atender múltiples conexiones.
Este programa tiene la finalidad de permitir la conexión de varios clientes, de
manera que se puedan comunicar entre ellos, todo esto siguiendo un protocolo
específico. Los saltos de línea son los que se tienen pensados que sean los que
delimiten el final de un mensaje. En la implementación de este proyecto se hizo
uso de la comunicación TCP y de las goroutines para atender conexiones concurrentemente.

\newpage
\section*{Introducción}
El principal problema que se busca resolver en este proyecto es poder tener un
programa que permita la comunicación entre dos o más clientes por medio de un chat. Esto se hará mediante un sistema separado de la vista del cliente que tiene
el objetivo de que el servidor se encargue de mantener los registros de clientes
que se encuentran conectados, mientras que el cliente envía solicitudes y
procesa cada mensaje recibido.

\noindent El alcance es local por el momento ya que solamente es posible establecer una
comunicación entre varios clientes dentro de una misma red Wi-Fi. En este caso,
ese sería una de las limitaciones que presenta nuestro programa, aunque no se
necesita hacer gran cosa para poder chatear entre distintos clientes.
\\[0.5cm]
\noindent Los objetivos principales se presentan a continuación:
\begin{itemize}
  \item Establecer una comunicación satisfacoria TCP entre varios clientes y el servidor
  \item Con base en el protocolo definir los mensajes JSON con respuestas y procesos
  \item Administrar clientes, sus estados y las salas que se crearán después
  \item Manejar posibles errores como de identificación o clientes inexistentes
  \item Organizar el código y asignar nombres coherentes a todo\\[0.1cm]
\end{itemize}

\noindent Además, se usará el MVC (Modelo-Vista-Controlador) para poder más fácilmente
organizar el trabajo que cada parte debe hacer. Cabe mencionar que este proyecto
se hará con el paradigma de Programación Orientado a Objetos (POO).

\newpage
\section*{Análisis y diseño}


\subsection*{Funcionalidades básicas}

\noindent Lo siguiente está dado por el protocolo.

\begin{enumerate}
  \item Identificar cada cliente con un nombre único máximo de 8 caracteres
  \item Consultar que clientes están conectados y ver sus estados
  \item Cambiar el estado de un cliente, con las opciones: \texttt{ACTIVE}, \texttt{AWAY} y \texttt{BUSY}
  \item Enviar mensajes públicos y privados
  \item Crear salas e invitar usuarios a unirse a salas, y pedir la lista de clientes
  \item Enviar mensajes a los clientes de una sala y/o abandonarla
  \item Notificar la desconexión, así también como la llegada de cada cliente
\end{enumerate}

\subsection*{Mensajes JSON}
Aquí tenemos una comunicación TCP, donde los mensajes están (deben estarlo)
codificados como un JSON válido que está delimitado por un salto de línea.
De nuevo, si la longitud del nombre de un usuario es mayor a 8 se harán las
operaciones indicadas por el protocolo, como es que si ocurre un error de
la falta de un argumento esperado después de un indicador de comando se debe
mandar un mensaje JSON y desconectar al cliente.

\subsection*{Estructura general}
La arquitectura se divide en el cliente y el servidor (adicionalmente
un paquete que se comparte entre ellos). En el cliente se interpretan los
mensajes recibidos por parte del usuario y se envían mensajes. El servidor es
la parte que se encarga de procesar múltiples conexiones, y una vez que se
valida la identificación del cliente procede a recibir los mensajes del cliente,
actualizar estados, mandar mensajes acorde a la petición del cliente, etc.\\
El paquete compartido define la estructura de los mensajes.

\begin{center}
\begin{tikzpicture}[node distance=13mm and 15mm]
\node[bloque] (cliente) {\textbf{Cliente}\\Entrada y salida};
\node[bloque,right=of cliente] (tcp) {Conexión TCP\\con JSON};
\node[bloque,right=of tcp] (servidor) {\textbf{Servidor}\\Validación y despacho};
\node[bloque,below=of servidor] (estado) {Estado del servidor\\Usuarios y salas};
\node[bloque,below=of tcp] (comun) {\texttt{\textbf{Paquete compartido}\\Estructura del Mensaje}};
\draw[flecha] (cliente) -- (tcp);
\draw[flecha] (tcp) -- (servidor);
\draw[flecha] (servidor) -- (estado);
\draw[flecha] (comun) -- (tcp);
\end{tikzpicture}
\end{center}

\newpage

\subsection*{Estructura del mensaje}
La estructura del \texttt{Mensaje} representa los campos del protocolo, como el
tipo de operación (\texttt{Operation}), el nombre de usuario (\texttt{Username}), el resultado (\texttt{Result}) y los datos adicionales (\texttt{Extra}).
El servidor básicamente maneja un map de \texttt{ClienteConectado} y otro de salas.
Cada uno de los clientes tiene su nombre, estado y conexión. Y cada sala mantiene su
nombre y los clientes que pertenecen a esa sala.

\begin{center}
\begin{tikzpicture}[node distance=10mm and 12mm]
\node[bloque] (server) {\textbf{Servidor}\\ - Puerto\\ - Clientes\\ - Salas};
\node[bloque,below left=of server] (cliente) {\textbf{ClienteC...}\\ - Nombre\\ - Estado\\ - Conexión};
\node[bloque,below right=of server] (sala) {\textbf{Sala}\\ - Nombre\\ - ClienteC...};
\draw[flecha] (server) -- node[left,align=center]{} (cliente);
\draw[flecha] (server) -- node[right,align=center]{} (sala);
\end{tikzpicture}
\end{center}

\subsection*{Proceso de identificación}
Cuando se establece una conexión, el cliente en automático envía una solicitud
de identificación. Luego de esto, el servidor comprueba si el JSON recibido es
válido y que lo primero que se haga sea identificarse con el tipo \texttt{IDENTIFY},
que el nombre no sea una cadena vacía, y también se cubre el caso cuando
el nombre con el que se intenta identificar se encuentra en uso por otro cliente.
En el caso de que todo salga bien, se registra al cliente y se responde con un mensaje
de éxito. De lo contrario, se envía el error correspondiente y se cierra la conexión
con el servidor.


\begin{center}
\begin{tikzpicture}[node distance=15mm]
\node[estado] (step1) {Aceptar conexión};
\node[estado,below=of step1] (identify) {Recibir el IDENTIFY};
\node[estado,below=of identify] (validar) {Validar mensaje y nombre\\JSON};
\node[estado,below left=of validar,xshift=-8mm] (error) {Enviar error\\y cerrar};
\node[estado,below right=of validar,xshift=8mm] (bien) {Registrar cliente\\y enviar SUCCESS};
\draw[flecha] (step1) -- (identify);
\draw[flecha] (identify) -- (validar);
\draw[flecha] (validar) -- node[above left]{inválido} (error);
\draw[flecha] (validar) -- node[above right]{válido} (bien);
\end{tikzpicture}
\end{center}

\subsection*{Comunicación cliente-servidor}
Se muestra en el siguiente esquema el intercambio de mensajes entre el servidor
y el cliente en el proceso de identificación, donde se sigue el algoritmo
de solicitar algo y obtener una respuesta dependiendo del mensaje.

\begin{center}
\begin{tikzpicture}[x=1.2cm,y=1.5cm]
\draw[dashed] (0,0) -- (0,-5.7);
\draw[dashed] (5,0) -- (5,-5.7);
\node[align=center] at (0,0.35) {\textbf{Cliente}};
\node[align=center] at (5,0.35) {\textbf{Servidor}};
\draw[flecha] (0,-1) -- node[above]{Conexión TCP} (5,-1);
\draw[flecha] (0,-2) -- node[above]{IDENTIFY} (5,-2);
\draw[flecha] (5,-3) -- node[above]{RESPONSE / SUCCESS} (0,-3);
\draw[flecha] (5,-4) -- node[above,align=center]{NEW\_USER\\(aviso)} (0,-4);
\node[align=center,text width=5cm] at (2.5,-5.1) {El cliente sigue enviando\\comandos y recibiendo mensajes};
\end{tikzpicture}
\end{center}

\subsection*{Tarjetas CRC (Clase-Responsabilidad-Colaborador)}

\noindent He aquí están las tarjetas CRC que nos permiten visualizar las
responsabilidades de cada una de las clases, así como las clases con las
que están relacionadas (colaborando).\\

\noindent \begin{tabularx}{\textwidth}{|>{\raggedright\arraybackslash}p{0.23\textwidth}|X|X|}
\hline
\textbf{Clase} & \textbf{Responsabilidades} & \textbf{Colaboradores}\\
\hline
Servidor & Aceptar conexiones. Validar mensajes. Administrar usuarios y también salas. Mandar mensajes & ClienteConectado, Sala, Mensaje\\
\hline
ClienteConectado & Mantener el nombre, estado y conexión del cliente & Servidor\\
\hline
Sala & Mantener a los integrantes, saber quiénes han sido invitados & Servidor y ClienteConectado\\
\hline
Cliente & Conectarse al servidor. Enviar comandos, recibir y mostrar mensajes & Servidor y Mensaje\\
\hline
Mensaje & Representar las solicitudes, respuestas y notificaciones dadas en el protocolo & Cliente y Servidor\\
\hline
\end{tabularx}

\newpage
\subsection*{Estados del cliente}
Como estado por default de un cliente (que ya se ha identificado) es \texttt{ACTIVE}.
Y puede cambiar a \texttt{AWAY} o \texttt{BUSY} por medio de una solicitud válida,
pues si se intenta solicitar por ejemplo un estado \texttt{BORED} que no está entre
las opciones de un estado se desconectará al cliente luego de dar el aviso de mensaje
inválido. Si se efectúa el cambio de estado, se debe de notificar a los demás clientes.

\begin{center}
\begin{tikzpicture}[node distance=30mm and 35mm]
\node[estado] (active) {ACTIVE};
\node[estado,right=of active] (away) {AWAY};
\node[estado,below=of away] (busy) {BUSY};
\draw[flecha,<->] (active) -- node[above]{STATUS} (away);
\draw[flecha,<->] (away) -- node[right]{STATUS} (busy);
\draw[flecha,<->] (busy) -- node[below left]{STATUS} (active);
\end{tikzpicture}
\end{center}

\newpage
\section*{Implementación}

\subsection*{Tecnologías usadas y organización}
El programa lo he desarrollado en el lenguaje de programación Go, donde he hecho
uso del paquete \texttt{net} para las conexiones TCP, del paquete \texttt{encoding/json}
para codificar y decodificar los mensajes, de las goroutines para poder atender
conexiones concurrentemente. En este proyecto se separa el código del servidor, del
cliente y de lo que se comparte entre ellos para tener una mejor visibilidad de lo
que está pasando en el programa. Lo siguiente es un esquema donde se puede apreciar
la atención de conexiones dentro del programa del servidor.

\begin{lstlisting}[style=gocode,caption={Esquema simplificado de atención de conexiones}]
listener, err := net.Listen("tcp", direccion)
if err != nil {
    return err
}
defer listener.Close()

for {
    conn, err := listener.Accept()
    if err != nil {
        continue
    }
    go servidor.AtiendeConexion(conn)
}
\end{lstlisting}

\subsection*{Manejo de mensajes}
Cuando se representa cada mensaje, se hace mediante una estructura compartida
transmitiéndose como JSON con al menos un campo. El campo \texttt{type} nos dice
la clase de mensaje recibido, el campo \texttt{Operation} nos dice la acción
que se intentó realizar, el campo \texttt{Username} nos dice el nombre de usuario, y así con los demás campos establecidos por el protocolo.
Lo que hacemos cuando uno o varios campos no son necesarios que aparezcan, es
desde el paquete compartido de Mensaje, etiquetarlos con \textbf{omitempty} y así
se omite cuando el campo es vacío.

\subsection*{Operaciones principales}
\noindent\begin{tabularx}{\textwidth}{|>{\ttfamily}p{0.30\textwidth}|X|}
\hline
\textnormal{\textbf{Operación}} & \textbf{Descripción}\\
\hline
IDENTIFY & Identifica a un cliente (automático al conectarse)\\
STATUS & Modifica el estado del cliente\\
USERS & Solicita la lista de clientes conectados\\
TEXT / PUBLIC\_TEXT & Envía mensajes privados / públicos\\
NEW\_ROOM & Crea una sala para chatear\\
INVITE / JOIN\_ROOM & Invita clientes a una sala / permite su ingreso a ella\\
ROOM\_USERS / ROOM\_TEXT & Consulta integrantes / envía mensajes a una sala\\
LEAVE\_ROOM / DISCONNECT & Deja una sala / Desconecta al cliente\\
\hline
\end{tabularx}

\begin{table}[ht]
    \centering
    \caption{Comandos y argumentos de los mensajes del protocolo}
    \label{tab:argumentos}
    \begin{tabular}{|l|p{6cm}|p{6cm}|}
        \hline
        \textbf{Operación} & \textbf{Uso} & \textbf{Ejemplo} \\
        \hline
        IDENTIFY & Automático al conectarse & \texttt{\{"type":"IDENTIFY","username":"Juan"\}} \\
        \hline
        STATUS & STATUS <estado> & STATUS BUSY \\
        \hline
        USERS & USERS & USERS \\
        \hline
        PUBLIC\_TEXT & <texto> & Hola a todos \\
        \hline
        TEXT & TEXT <username> <mensaje> & TEXT Kay Hola, ¿cómo estás? \\
        \hline
        NEW\_ROOM & NEW\_ROOM <roomname> & NEW\_ROOM Tareas \\
        \hline
        INVITE & INVITE <roomname> <usernames> & INVITE Tareas Kay \\
        \hline
        JOIN\_ROOM & JOIN\_ROOM <roomname> & JOIN\_ROOM Tareas \\
        \hline
        ROOM\_USERS & ROOM\_USERS <roomname> & ROOM\_USERS Tareas \\
        \hline
        ROOM\_TEXT & ROOM\_TEXT <roomname> <mensaje> & ROOM\_TEXT Tareas Me pasas la tarea? \\
        \hline
        LEAVE\_ROOM & LEAVE\_ROOM <roomname> & LEAVE\_ROOM Tareas \\
        \hline
        DISCONNECT & DISCONNECT & DISCONNECT \\
        \hline
    \end{tabular}
\end{table}

\subsection*{Validación y manejo de posibles errores}
El servidor rechaza mensajes JSON inválidos (mal construidos), solicitudes
distintas de la identificación como primer paso, nombres vacíos o repetidos
y operaciones con parámetros incorrectos. También se comprueba la existencia
de los clientes a los que se les efectúa una determinada acción, como invitar
a una sala, así como la comprobación de la pertenencia del cliente a la sala
a la que quiere invitar a otro cliente. Así con los casos que lo requieran.

\subsection*{Requerimientos para el programa}
Solamente ejecutar los dos archivos \texttt{Servidor} y \texttt{Cliente}, pero
comenzando con el del servidor. Se debe hacer \texttt{./Servidor} y luego con
el del cliente \texttt{./Cliente IP:HOST <nombre>}, suponiendo que están
conectados a la misma red Wi-Fi. Solamente recordando que el nombre debe ser a
lo más de 8 caracteres de longitud, y dentro del programa, al crear una sala
su nombre debe ser máximo de 16 caracteres.


\newpage
\section*{Conclusiones}
Después del desarrollo de este proyecto, aprendí con más claridad cómo es
que funciona la comunicación TCP, en este caso del cliente-servidor. Además,
pude aprender (no en su totalidad) como programar en Go, y una técnica usando
net.Pipe() en go, que es básicamente simular una conexión entre dos extremos en
memoria, esto me ayudó a realizar pequeñas pruebas sin necesidad de establecer
una conexión real. Lo que se me dificultó fue entender como era esa
comunicación que al final de cuentas era lo principal en este proyecto.\\

\noindent Veo un área de mejora en mis comentarios dentro del proyecto, y también en
como nombro tanto las variables como las funciones, además de organizar mis
ideas. Si en algún futuro pudiera hacer de este proyecto algo mejor, lo que
haría sería darle más alcance y ponerle al menos una interfaz bonita.
Lo que me ayudó a visualizar el funcionamiento de este proyecto fue usar el
paradigma de POO y el MVC, así también como los distintos diagramas y tarjetas
CRC. Otra cosa que se me dificultó fue cubrir todos los posibles casos dados
implícita y explícitamente en el protocolo, pues al cubrir unos salían otros
nuevos que no había contemplado.



\end{document}
